\idxletter{A}
\sourcepage{744}
\idxentry{0}{先验假设}{A priori hypothesis}{54, 395g}
\idxentry{0}{绝对风险}{Absolute risk}{250}
\idxentry{1}{降低}{reduction}{132, 395g}
\idxentry{0}{可接触人群}{Accessible population}{27, 395g}
\idxentry{0}{准确度／准确率}{Accuracy}{9, 39, 43–45, 44f, 44t, 240, 395g}
\idxentry{1}{医学检验的准确度}{of medical tests}{}
\idxentry{2}{绝对风险、风险比、风险差与风险率比}{absolute risks, risk ratios, risk differences, and hazard ratios}{250}
\idxentry{2}{分析}{analysis}{247–250}
\idxentry{2}{校准}{calibration}{250}
\idxentry{2}{设计}{designs}{245–247}
\idxentry{2}{似然比}{likelihood ratios}{249–250, 249t}
\idxentry{2}{结局变量}{outcome variable}{246–247}
\idxentry{2}{预测变量}{predictor variable}{246}
\idxentry{2}{受试者工作特征曲线}{receiver operating characteristic curves}{248, 248f}
\idxentry{2}{抽样}{sampling}{245–246}
\idxentry{2}{灵敏度、特异度、阳性与阴性预测值及准确率}{sensitivity, specificity, positive and negative predictive values, and accuracy}{247, 247t}
\idxentry{1}{提高准确度的策略}{strategies for enhancing}{45–46, 45t}
\idxentry{0}{ACE2。参见血管紧张素转换酶2}{ACE2. See Angiotensin-converting enzyme 2 (ACE2)}{}
\idxentry{0}{活性对照试验}{Active control trial}{221–224, 223f, 395g}
\idxentry{0}{ACTT-1。参见适应性COVID-19治疗试验}{ACTT-1. See Adaptive COVID-19 Treatment Trial (ACTT-1)}{}
\idxentry{0}{适应性临床试验设计}{Adaptive clinical trial design}{224–225, 395g}
\idxentry{0}{适应性COVID-19治疗试验（ACTT-1）}{Adaptive COVID-19 Treatment Trial (ACTT-1)}{224}
\idxentry{0}{适应性随机化}{Adaptive randomization}{204, 395g}
\idxentry{0}{加法与乘法交互作用}{Additive vs. multiplicative interactions}{177}
\idxentry{0}{方案依从性}{Adherence to protocol}{207–209, 207t}
\idxentry{0}{调整}{Adjustment}{74, 185, 395g}
\idxentry{0}{自报测量的实施}{Administering self-reported measures}{}
\idxentry{1}{实施平台}{administration platforms}{319}
\idxentry{1}{访谈方式}{interviewing approaches}{320}
\idxentry{1}{自填与访谈员施测}{self- vs. interviewer administration}{319}
\idxentry{0}{实施平台}{Administration platforms}{319}
\idxentry{0}{研究申请书的行政部分}{Administrative parts, of research proposal}{356–358}
\sourcepage{745}
\idxentry{0}{事先提出的假设}{Advance hypothesis}{54}
\idxentry{0}{不良事件}{Adverse events}{199–200, 396g}
\idxentry{0}{事后提出的假设}{After-the-fact hypothesis}{54}
\idxentry{0}{AI。参见人工智能}{AI. See Artificial intelligence (AI)}{}
\idxentry{0}{All of Us研究计划}{All of Us Research Program}{295}
\idxentry{0}{\ensuremath{\alpha}水平}{Alpha (\ensuremath{\alpha})}{56–57, 57t, 65, 396g}
\idxentry{0}{其他干预性研究设计}{Alternate interventional study designs}{}
\idxentry{1}{贝叶斯临床试验}{Bayesian clinical trials}{235–238, 236t, 237f}
\idxentry{1}{非随机干预设计。参见非随机干预设计}{nonrandom interventional designs. See Nonrandom interventional designs}{}
\idxentry{0}{备择假设}{Alternative hypothesis}{55, 64, 65, 396g}
\idxentry{1}{假设的侧数}{sides of}{58}
\idxentry{0}{其他测量策略}{Alternative measurement strategies}{320}
\idxentry{0}{其他随机设计}{Alternative randomized designs}{}
\idxentry{1}{活性对照试验}{active control trial}{221–224, 223f}
\idxentry{1}{适应性设计}{adaptive designs}{224–225}
\idxentry{1}{整群随机设计}{cluster randomized designs}{227–228, 227f}
\idxentry{1}{交叉设计}{crossover designs}{225, 225f}
\idxentry{1}{等效性}{equivalence}{221–224, 223f}
\idxentry{1}{析因试验}{factorial trial}{221, 222f}
\idxentry{1}{单例设计}{N-of-1 designs}{226}
\idxentry{1}{非劣效性}{noninferiority}{221–224, 223f}
\idxentry{1}{阶梯楔形设计}{stepped wedge designs}{228–229, 228f}
\idxentry{1}{等待名单设计}{wait-list designs}{226}
\idxentry{0}{阿尔茨海默病神经影像倡议}{Alzheimer’s Disease Neuroimaging Initiative}{296}
\idxentry{0}{分析阶段}{Analysis phase}{}
\idxentry{1}{混杂因素的处理}{coping with confounders in}{172t}
\idxentry{2}{多变量建模}{multivariable modeling}{173}
\idxentry{2}{倾向评分}{propensity scores}{173–175}
\idxentry{2}{分层}{stratification}{171–173}
\idxentry{1}{减少偏倚所致错误}{minimizing errors due to bias}{163}
\idxentry{0}{分析性研究}{Analytic study}{5, 53, 116, 396g}
\idxentry{1}{样本量估计方法}{sample size techniques for}{65, 66t}
\idxentry{2}{卡方检验}{chi-squared test}{68–71, 87–89}
\idxentry{2}{相关系数}{correlation coefficient}{71–72, 90}
\idxentry{2}{t检验}{t test}{66–68, 86, 93}
\idxentry{0}{附属研究}{Ancillary study}{291, 301–302, 396g}
\idxentry{0}{血管紧张素转换酶2（ACE2）}{Angiotensin-converting enzyme 2 (ACE2)}{148–149}
\idxentry{0}{应用程序编程接口（API）}{Application programming interface (API)}{300}
\idxentry{0}{研究策略中的研究方法部分}{Approach section, of research strategy}{358–359}
\idxentry{0}{人工智能（AI）}{Artificial intelligence (AI)}{110}
\idxentry{0}{按实际治疗分析}{As-treated analysis}{211–212, 213f, 396g}
\sourcepage{746}
\idxentry{0}{关联}{Association}{5, 157, 396g}
\idxentry{1}{度量指标}{measures of}{132}
\idxentry{1}{观察性研究关联与因果效应的区别}{in observational study vs. causal effect}{158–159, 159t}
\idxentry{1}{关联强度}{strength of}{178}
\idxentry{0}{署名}{Authorship}{106–107}
\idxletter{B}
\idxentry{0}{保存的标本}{Banked specimens}{202}
\idxentry{0}{基线变量}{Baseline variables}{}
\idxentry{1}{随机盲法试验}{randomized blinded trial}{}
\idxentry{2}{额外测量}{additional measurements}{202}
\idxentry{2}{保存的标本}{banked specimens}{202}
\idxentry{2}{参与者}{participants}{202}
\idxentry{0}{贝叶斯分析}{Bayesian analysis}{224, 396g}
\idxentry{0}{贝叶斯方法}{“Bayesian” approach}{61}
\idxentry{0}{贝叶斯试验}{Bayesian trial}{211, 224, 235–238, 236t, 237f, 396g}
\idxentry{0}{前后试验设计}{Before-after (or pre-post) trial designs}{230–231, 230f}
\idxentry{0}{前后或中断时间序列设计}{Before-after or interrupted time series designs}{206}
\idxentry{0}{前后研究（试验）设计}{Before-after study (or trial) design}{}
\idxentry{1}{临床决策的前后研究}{of clinical decision-making}{253–254}
\idxentry{1}{非等效对照组。参见受控前后研究}{with non-equivalent control group. See Controlled before-after study}{}
\idxentry{0}{有利原则}{Beneficence}{96, 396–397g}
\idxentry{0}{\ensuremath{\beta}水平}{Beta (\ensuremath{\beta})}{56–57, 57t, 65, 397g}
\idxentry{0}{组间设计}{Between-group design}{225}
\idxentry{0}{偏倚}{Bias}{9, 56, 158, 159t, 397g}
\idxentry{1}{尽量减少所致错误}{minimizing errors due to}{161f}
\idxentry{2}{分析阶段}{analysis phase}{163}
\idxentry{2}{设计阶段}{design phase}{161–163}
\idxentry{0}{大数据}{Big data}{110}
\idxentry{0}{生物样本库研究}{Biobanking research}{265}
\idxentry{0}{生物学合理性}{Biologic plausibility}{179}
\idxentry{0}{生物标志物}{Biomarker}{4, 198, 397g}
\idxentry{0}{Bland–Altman图}{Bland-Altman plot}{244}
\idxentry{0}{盲法}{Blinding}{45t, 46, 144, 196, 397g}
\idxentry{1}{病例对照研究中的盲法}{in case-control study}{144t}
\idxentry{1}{难以或无法设盲}{difficult or impossible}{205–206}
\idxentry{1}{医学检查的准确度}{of medical tests}{242}
\idxentry{1}{实验室程序的质量控制}{quality control for laboratory procedures}{333}
\idxentry{1}{随机盲法试验}{randomized blinded trial}{204–206}
\idxentry{0}{区组随机化}{Blocked randomization}{203, 397g}
\idxentry{0}{Bonferroni校正}{Bonferroni correction}{61, 397g}
\idxentry{0}{自助法}{Bootstrapping}{397g}
\sourcepage{747}
\idxentry{0}{分支问题}{Branching questions}{309–310}
\idxentry{0}{预算}{Budgets}{353}
\idxentry{1}{申请书的行政部分}{administrative parts of proposal}{356–357}
\idxentry{1}{预算理由}{justification}{357, 358}
\idxentry{1}{预算编制}{preparation of}{357}
\idxletter{C}
\idxentry{0}{校准}{Calibration}{42, 45t, 46, 245, 250, 397g}
\idxentry{0}{癌基因中的多态性}{Cancer genes, polymorphisms in}{189}
\idxentry{0}{癌症登记}{Cancer registries}{294–295}
\idxentry{0}{心脏代谢疾病}{Cardiometabolic diseases}{148–149}
\idxentry{0}{心血管疾病登记}{Cardiovascular disease registries}{295}
\idxentry{0}{职业发展K类资助申请的评分部分}{Career development “K” grant application scored sections}{363t}
\idxentry{0}{残留效应}{Carryover effects}{225, 398g}
\idxentry{0}{CART分析。参见分类与回归树分析}{CART analysis. See Classification and regression tree (CART) analysis}{}
\idxentry{0}{病例}{Case}{4, 398g}
\idxentry{1}{病例寻找与抽样}{finding and sampling}{141–142, 141f}
\idxentry{0}{病例队列研究}{Case-cohort study}{148, 398g}
\idxentry{0}{病例对照研究}{Case-control study}{4, 4t, 128, 138, 398g}
\idxentry{1}{设盲方法}{approaches to blinding in}{144t}
\idxentry{1}{基础}{basics}{138–140, 139f}
\idxentry{1}{病例交叉研究}{case-crossover studies}{149–150}
\idxentry{1}{其中的卡方检验}{chi-squared test in}{69}
\idxentry{1}{习题}{exercises for}{156}
\idxentry{1}{发病密度巢式研究}{incidence-density nested}{146–148, 147f}
\idxentry{2}{优点}{strengths}{149}
\idxentry{2}{缺点}{weaknesses}{149}
\idxentry{1}{匹配}{matching}{166}
\idxentry{1}{每个病例配多个对照}{multiple controls per case in}{81}
\idxentry{1}{巢式研究}{nested}{145, 146f}
\idxentry{2}{优点}{strengths}{149}
\idxentry{2}{缺点}{weaknesses}{149}
\idxentry{1}{巢式病例队列设计}{nested case-cohort design}{148–149}
\idxentry{2}{优点}{strengths}{149}
\idxentry{2}{缺点}{weaknesses}{149}
\idxentry{1}{观察性设计}{observational designs}{}
\idxentry{2}{优缺点}{advantages and disadvantages of}{151–152t}
\idxentry{2}{设计的选择}{choosing among}{150}
\idxentry{1}{比值比}{odds ratios for}{153–154}
\idxentry{2}{与风险比的关系}{risk ratio and}{155}
\idxentry{1}{优点}{strengths of}{140}
\idxentry{2}{产生假设的用途}{generating hypothesis, usefulness for}{140}
\sourcepage{748}
\idxentry{2}{研究罕见结局的效率}{rare outcomes, efficiency for}{140}
\idxentry{1}{t检验}{t test}{67}
\idxentry{1}{缺点}{weaknesses of}{}
\idxentry{2}{差异性测量偏倚}{differential measurement bias}{143–145}
\idxentry{2}{选择偏倚}{selection bias}{141–143, 141f}
\idxentry{0}{病例交叉研究}{Case-crossover study}{138, 149–150, 398g}
\idxentry{0}{病例报告表（CRF）}{Case report form (CRF)}{345}
\idxentry{0}{病例系列}{Case series}{138, 398g}
\idxentry{0}{分类变量}{Categorical variable}{40, 72–73, 398g}
\idxentry{1}{分类变量的医学检查}{medical tests of}{244}
\idxentry{0}{因果效应}{Causal effect}{157, 398g}
\idxentry{1}{与观察性研究关联的区别}{association in observational study vs.}{158–159, 159t}
\idxentry{1}{混杂}{confounding}{164}
\idxentry{1}{混杂与效应修饰}{confounding and effect modification}{191–192}
\idxentry{1}{分析阶段处理混杂}{coping with confounders in analysis phase}{172t}
\idxentry{2}{多变量建模}{multivariable modeling}{173}
\idxentry{2}{倾向评分}{propensity scores}{173–175}
\idxentry{2}{分层}{stratification}{171–173}
\idxentry{1}{设计阶段处理混杂}{coping with confounders in design phase}{164–165, 165t}
\idxentry{2}{匹配}{matching}{166–167}
\idxentry{2}{机会性研究}{opportunistic studies}{168–171}
\idxentry{2}{限定}{specification}{165–166}
\idxentry{1}{理解因果关系的反事实框架}{counterfactual framework for understanding causality}{157–158, 158f}
\idxentry{1}{有向无环图}{directed acyclic graphs}{}
\idxentry{2}{混杂}{confounding}{185, 185f}
\idxentry{2}{连接变量}{connecting variables}{184–186, 184–185f}
\idxentry{2}{循环图}{cyclic graph}{185, 185f}
\idxentry{2}{例子}{example}{182, 188–190}
\idxentry{2}{因果推断}{inferences about causality}{186–187, 186–187t}
\idxentry{2}{局限}{limitations}{190}
\idxentry{2}{中介作用}{mediation}{184, 184f}
\idxentry{2}{未连接全部变量}{not connecting variables}{182–184, 183–184f}
\idxentry{2}{与抽样的关系}{sampling and}{187}
\idxentry{2}{共同效应}{shared effects}{184–185, 185f}
\idxentry{2}{共有原因}{sharing cause}{188}
\idxentry{1}{果因关系}{effect-cause}{163–164, 163t}
\idxentry{1}{效应修饰与混杂}{effect modification, confounding and}{191–192}
\idxentry{1}{尽量减少错误}{errors, minimizing}{}
\idxentry{2}{偏倚所致错误}{due to bias}{161–163, 161f}
\idxentry{2}{偶然所致错误}{due to chance}{160, 160t}
\idxentry{1}{支持因果关系的依据}{evidence favoring causality}{178–179}
\sourcepage{749}
\idxentry{1}{习题}{exercises for}{195}
\idxentry{1}{辨明因果方向}{getting direction of causality right}{163–164, 163t}
\idxentry{1}{建模的简化例子}{modeling, simplified example of}{193–194}
\idxentry{1}{量化的陷阱}{pitfalls in quantifying}{}
\idxentry{2}{对共同效应条件化}{conditioning on shared effect}{175–176, 175f}
\idxentry{2}{效应修饰与交互作用}{effect modification and interaction}{176}
\idxentry{2}{低估因果效应}{underestimation of causal effects}{176–177}
\idxentry{1}{策略选择}{strategy, choosing}{178}
\idxentry{0}{因果推断}{Causal inference}{8}
\idxentry{0}{因果关系}{Causality}{}
\idxentry{1}{理解因果关系的反事实框架}{counterfactual framework for understanding}{157–158, 158f}
\idxentry{1}{支持因果的依据}{evidence favoring}{178–179}
\idxentry{1}{辨明方向}{getting direction of right}{163–164, 163t}
\idxentry{1}{有关推断}{inferences about}{186–187}
\idxentry{0}{因果关系}{Cause-effect}{5, 398g}
\idxentry{0}{CDC WONDER。参见美国疾病控制与预防中心流行病学研究广域在线数据}{CDC WONDER. See Centers for Disease Control and Prevention’s Wide-ranging Online Data for Epidemiologic Research (CDC WONDER)}{}
\idxentry{0}{删失}{Censoring}{123, 131, 398g}
\idxentry{0}{美国疾病控制与预防中心流行病学研究广域在线数据（CDC WONDER）}{Centers for Disease Control and Prevention’s Wide-ranging Online Data for Epidemiologic Research (CDC WONDER)}{294}
\idxentry{0}{保密证书}{Certificate of confidentiality}{102, 398g}
\idxentry{0}{偶然}{Chance}{158, 159t}
\idxentry{1}{尽量减少所致错误}{minimizing errors due to}{160, 160t}
\idxentry{0}{卡方检验}{Chi-squared test}{68–71, 87–89, 399g}
\idxentry{0}{儿童研究}{Children, research on}{102–103}
\idxentry{0}{对照选择}{Choice of control}{197–198}
\idxentry{0}{干预选择}{Choice of intervention}{196–197}
\idxentry{0}{结局选择}{Choice of outcome}{255}
\idxentry{0}{问题的清晰程度}{Clarity, of questions}{310}
\idxentry{0}{分类与回归树}{Classification and regression tree (CART)}{251。另见递归分割（Recursive partitioning）}
\idxentry{0}{诊所来源对照}{Clinic-based control}{142, 399g}
\idxentry{0}{临床数据管理员}{Clinical data managers}{340, 399g}
\idxentry{0}{临床决策}{Clinical decisions}{}
\idxentry{1}{检查结果对临床决策作用的研究}{studies of effect of test results on}{}
\idxentry{2}{临床决策的前后研究}{before/after studies of clinical decision-making}{253–254}
\idxentry{2}{诊断收益研究}{diagnostic yield studies}{253}
\idxentry{0}{临床经验}{Clinical experience}{19}
\idxentry{0}{临床监察员}{Clinical monitors}{336}
\idxentry{0}{临床预测模型}{Clinical prediction model}{251, 399g}
\idxentry{1}{建立模型}{developing}{251, 252f}
\sourcepage{750}
\idxentry{1}{验证}{validation}{252–253}
\idxentry{0}{临床程序质量控制}{Clinical procedures, quality control for}{332–333, 333t}
\idxentry{0}{临床研究。另见研究}{Clinical research. See also Research}{}
\idxentry{1}{定性方法。参见定性方法}{qualitative approaches in. See Qualitative approaches}{}
\idxentry{1}{相关法规}{regulations on}{95}
\idxentry{1}{实施所需资源}{resources for conducting}{325}
\idxentry{0}{合同研究组织（CRO）}{Clinical research organization (CRO)}{366}
\idxentry{0}{从临床到人群研究}{Clinical to population research}{11}
\idxentry{0}{临床试验}{Clinical trial}{3, 399g}
\idxentry{1}{随访与方案依从性}{follow-up and adherence to the protocol}{207–209, 207t}
\idxentry{1}{嵌入医疗服务}{in health care}{206}
\idxentry{1}{意向治疗、符合方案和按实际治疗分析}{intention-to-treat, per-protocol, and as-treated analysis}{211–212, 212–213f}
\idxentry{1}{监测}{monitoring}{209–211}
\idxentry{1}{传统临床机构之外}{outside of traditional clinical sites}{209}
\idxentry{1}{排除参与者的理由}{reasons for excluding people from}{201t}
\idxentry{1}{选择标准}{selection criteria for}{200–201}
\idxentry{1}{检查对结局作用的研究}{studies of effect of testing on outcomes}{256}
\idxentry{1}{亚组分析}{subgroup analysis}{213–214}
\idxentry{0}{临床人群与有代表性的人群}{Clinical vs. representative populations}{30–31}
\idxentry{0}{结束阶段}{Close-out}{330–331, 399g}
\idxentry{0}{封闭式问题}{Closed-ended question}{269, 308, 399g}
\idxentry{1}{有序回答选项}{ordered response options for}{308–309}
\idxentry{1}{回答选项}{response options}{308}
\idxentry{0}{整群随机化}{Cluster randomization}{227–228, 227f, 399g}
\idxentry{0}{整群抽样}{Cluster sampling}{32, 73, 399g}
\idxentry{0}{编码手册}{Codebook}{275, 399g}
\idxentry{0}{编码约定}{Coding conventions}{343–344, 344f}
\idxentry{0}{变异系数}{Coefficient of variation (CV)}{245, 399–400g}
\idxentry{0}{队列研究}{Cohort study}{4, 4t, 116, 132, 400g}
\idxentry{1}{分析方法}{approach to analyzing}{}
\idxentry{2}{关联指标}{association, measures of}{132}
\idxentry{2}{结局频率指标}{outcome frequency, measures of}{131–132}
\idxentry{1}{其中的卡方检验}{chi-squared test in}{69}
\idxentry{1}{数据表视图中的数据表}{data table in “datasheet view” for}{341f}
\idxentry{1}{双队列研究}{double-cohort studies}{126f, 127–128}
\idxentry{1}{习题}{exercises for}{137}
\idxentry{1}{外部对照}{external controls}{127–128}
\idxentry{1}{尽可能完成随访}{follow-up, maximizing}{132–133, 133t}
\idxentry{1}{多队列研究}{multiple-cohort studies}{127–128}
\idxentry{1}{前瞻性与回顾性研究}{prospective vs. retrospective studies}{119–122, 121t}
\idxentry{1}{复杂程度范围}{range of complexity of}{122–126, 123t, 124–125f}
\sourcepage{751}
\idxentry{1}{优缺点}{strengths and weaknesses of}{130–131}
\idxentry{1}{其中的t检验}{t test in}{66–67}
\idxentry{1}{治疗或干预的队列研究}{of treatments or interventions}{128–130}
\idxentry{0}{共同干预}{Cointervention}{196, 198, 400g}
\idxentry{0}{合作}{Collaborations}{19}
\idxentry{0}{多中心合作研究}{Collaborative multicenter studies}{336}
\idxentry{0}{合作研究连续谱}{Collaborative research continuum}{285–287, 285f, 286t}
\idxentry{0}{问卷中穷尽全部可能的回答}{Collectively exhaustive responses in questionnaire}{345}
\idxentry{0}{碰撞点}{Collider}{175。另见对共同效应条件化（Conditioning on a shared effect）}
\idxentry{1}{有向无环图}{directed acyclic graphs}{184}
\idxentry{1}{分层偏倚}{stratification bias}{175–176, 175f}
\idxentry{0}{结直肠癌}{Colorectal cancer}{159t, 189}
\idxentry{0}{共同数据模型}{Common data models}{298, 400g}
\idxentry{0}{社区咨询委员会}{Community advisory board}{285, 400g}
\idxentry{0}{基于社区的参与式研究（CBPR）}{Community-based participatory research (CBPR)}{286, 400g}
\idxentry{0}{社区对照组}{“Community” control group}{143}
\idxentry{0}{社区参与式研究}{Community-engaged research}{267, 282, 401g}
\idxentry{1}{研究方法}{approaches to}{284–285}
\idxentry{2}{合作研究连续谱}{continuum of collaborative research}{285–287, 285f, 286t}
\idxentry{1}{面临的挑战}{challenges of}{}
\idxentry{2}{建立信任}{building trust}{287}
\idxentry{2}{管理时间进度与期望}{managing timelines and expectations}{287–288}
\idxentry{1}{习题}{exercise for}{290}
\idxentry{1}{开展的理由}{reasons for}{}
\idxentry{2}{可推广性}{generalizability}{283}
\idxentry{2}{促进健康公平}{health equity, promoting}{284}
\idxentry{2}{地方能力与项目可持续性}{local capacity and program sustainability}{283–284}
\idxentry{2}{地方问题与地方知识}{local questions and local knowledge}{282–283, 283t}
\idxentry{0}{社区参与工作坊}{Community engagement studio}{285, 400g}
\idxentry{0}{社区外展}{Community outreach}{285, 400g}
\idxentry{0}{比较疗效试验}{Comparative efficacy trial}{221。另见活性对照试验（Active control trial）}
\idxentry{0}{比较性研究}{Comparative study}{268, 268t, 401g}
\idxentry{1}{指南}{guides}{273}
\idxentry{1}{撰写与分享结果}{writing and sharing results}{277}
\idxentry{0}{复杂队列研究}{Complex cohort studies}{124–125f}
\idxentry{1}{与简单队列的比较}{simple cohort study vs.}{123t}
\idxentry{0}{复杂假设}{Complex hypothesis}{53, 401g}
\idxentry{0}{队列研究复杂程度}{Complexity of cohort studies}{122–126, 123t}
\idxentry{0}{复合结局}{Composite outcome}{199, 401g}
\idxentry{0}{概念适当性}{Conceptual adequacy}{318, 401g}
\idxentry{0}{概念等价性}{Conceptual equivalence}{318, 401g}
\sourcepage{752}
\idxentry{0}{一致率}{Concordance}{167, 401g}
\idxentry{0}{条件化}{Conditioning}{175, 401g}
\idxentry{0}{对共同结果条件化}{Conditioning on a shared effect}{175–176, 175f, 401g}
\idxentry{0}{会议}{Conferences}{18}
\idxentry{0}{置信区间}{Confidence interval}{6, 75, 160, 401–402g}
\idxentry{0}{保密}{Confidentiality}{101–102}
\idxentry{1}{数据管理}{data management}{349}
\idxentry{1}{保护措施}{protection}{272–273}
\idxentry{0}{利益冲突}{Conflicts of interest}{95, 402g}
\idxentry{1}{财务方面}{financial}{107}
\idxentry{1}{专业方面}{professional}{107}
\idxentry{1}{应对}{responding to}{107–108}
\idxentry{1}{类型}{types of}{107}
\idxentry{0}{混杂因素}{Confounder}{74。另见混杂（Confounding）}
\idxentry{1}{分析阶段的处理}{coping with, in analysis phase}{172t}
\idxentry{2}{多变量建模}{multivariable modeling}{173}
\idxentry{2}{倾向评分}{propensity scores}{173–175}
\idxentry{2}{分层}{stratification}{171–173}
\idxentry{1}{设计阶段的处理}{coping with, in design phase}{164–165, 165t}
\idxentry{2}{匹配}{matching}{166–167}
\idxentry{2}{机会性研究}{opportunistic studies}{168–171}
\idxentry{2}{限定}{specification}{165–166}
\idxentry{0}{混杂}{Confounding}{158, 159t, 402g}
\idxentry{1}{有向无环图}{directed acyclic graphs}{185, 185f}
\idxentry{1}{与效应修饰的关系}{and effect modification}{191–192}
\idxentry{1}{演示用假想例子}{hypothetical examples to demonstrate}{191–192}
\idxentry{0}{适应证混杂}{Confounding by indication}{173–174, 297, 402g}
\idxentry{0}{混杂变量}{Confounding variables}{6, 196。另见混杂（Confounding）}
\idxentry{0}{连续样本}{Consecutive sample}{33, 402g}
\idxentry{0}{知情同意书}{Consent forms}{99}
\idxentry{0}{一致性}{Consistency}{41}
\idxentry{0}{构念效度}{Construct validity}{46, 317t, 402g}
\idxentry{0}{建构主义}{Constructivism}{266}
\idxentry{0}{沾染}{Contamination}{227, 402g}
\idxentry{0}{内容效度}{Content validity}{46, 317t, 402g}
\idxentry{0}{连续结局变量}{Continuous outcome variables}{66}
\idxentry{0}{连续预测变量}{Continuous predictor variable}{66}
\idxentry{0}{连续变量}{Continuous variable}{40, 76, 91, 402–403g}
\idxentry{1}{与二分类变量的比较}{dichotomous variables vs.}{78–79}
\idxentry{1}{分类变量的医学检查}{medical tests of}{244–245}
\idxentry{1}{减少样本量并提高效能}{for minimizing sample size and maximizing power}{78–79}
\sourcepage{753}
\idxentry{1}{连续变量描述性研究的样本量}{sample size for descriptive study of}{91}
\idxentry{0}{合作研究连续谱}{Continuum of collaborative research}{285–287, 285f, 286t}
\idxentry{0}{合同与美国国立卫生研究院}{Contracts, and National Institutes of Health}{362–364, 363t, 364f}
\idxentry{0}{对照条件}{Control conditions}{196}
\idxentry{1}{与干预的关系}{intervention and}{197f}
\idxentry{2}{对照选择}{choice of control}{197–198}
\idxentry{2}{干预选择}{choice of intervention}{196–197}
\idxentry{2}{共同干预}{cointerventions}{198}
\idxentry{0}{对照组的选择}{Control group, selection of}{143}
\idxentry{0}{有对照的前后研究}{Controlled before-after study}{232, 403g}
\idxentry{0}{对照；控制}{Control}{4, 403g}
\idxentry{1}{抽样}{sampling}{142–143}
\idxentry{0}{便利样本}{Convenience sample}{33, 403g}
\idxentry{0}{协调中心}{Coordinating center}{336, 403g}
\idxentry{0}{混杂因素的处理}{Coping with confounders}{}
\idxentry{1}{分析阶段}{in analysis phase}{172t}
\idxentry{2}{多变量建模}{multivariable modeling}{173}
\idxentry{2}{倾向评分}{propensity scores}{173–175}
\idxentry{2}{分层}{stratification}{171–173}
\idxentry{1}{设计阶段}{in design phase}{164–165, 165t}
\idxentry{2}{匹配}{matching}{166–167}
\idxentry{2}{机会性研究}{opportunistic studies}{168–171}
\idxentry{2}{限定}{specification}{165–166}
\idxentry{0}{相关系数}{Correlation coefficient}{71–72, 90, 403g}
\idxentry{0}{成本研究}{Costs, study of}{254}
\idxentry{0}{计数变量}{Count variable}{72–73, 403g}
\idxentry{0}{反事实框架}{Counterfactual framework}{157–158, 158f, 403g}
\idxentry{0}{协变量}{Covariate}{39, 116, 403g}
\idxentry{0}{覆盖分析}{Coverage analysis}{327}
\idxentry{0}{Cox模型}{Cox model}{74, 132, 403g}
\idxentry{0}{创造力}{Creativity}{19}
\idxentry{0}{可信区间}{Credible interval}{236, 404g}
\idxentry{0}{CRF。参见病例报告表}{CRF. See Case report form (CRF)}{}
\idxentry{0}{效标关联效度}{Criterion-related validity}{46, 404g}
\idxentry{0}{横断面研究}{Cross-sectional study}{4, 4t, 116, 404g}
\idxentry{1}{基本设计}{basic design}{116–118, 117f, 117t}
\idxentry{1}{其中的相关系数}{correlation coefficient in}{72}
\idxentry{1}{习题}{exercises for}{137}
\idxentry{1}{连续调查}{serial surveys}{119}
\idxentry{1}{优缺点}{strengths and weaknesses of}{118–119}
\idxentry{0}{交叉转组}{Crossover}{211, 404g}
\sourcepage{754}
\idxentry{0}{交叉研究}{Crossover study}{225, 225f, 404g}
\idxentry{0}{累积发病率。参见发病}{Cumulative incidence. See Incidence}{}
\idxentry{0}{CV。参见变异系数}{CV. See Coefficient of variation (CV)}{}
\idxentry{0}{循环图}{Cyclic graph}{185, 185f}
\idxletter{D}
\idxentry{0}{DAG。参见有向无环图}{DAG. See Directed Acyclic Graph (DAG)}{}
\idxentry{0}{数据}{Data}{}
\idxentry{1}{数据分析}{analysis of}{349}
\idxentry{1}{资料收集}{collection}{274t}
\idxentry{2}{保密保护}{confidentiality protection}{272–273}
\idxentry{2}{研究现场与参与者}{field sites and study participants}{272}
\idxentry{2}{访谈提纲}{interview guides}{273}
\idxentry{2}{访谈／焦点小组}{interviews/focus groups}{272}
\idxentry{2}{样本量的确定}{sample size, determination of}{271–272}
\idxentry{1}{多来源数据合并}{combining from multiple sources}{302}
\idxentry{1}{数据元素}{data elements}{344–345}
\idxentry{1}{数据编辑}{editing}{349}
\idxentry{1}{数据提取}{extracting}{348f}
\idxentry{2}{识别并纠正数据错误}{identifying and correcting errors in data}{348–349}
\idxentry{1}{管理、分析与报告}{management, analysis, and reporting}{}
\idxentry{2}{管理与分析过程}{process of managing and analyzing}{276–277}
\idxentry{2}{定性分析目标}{qualitative analysis, goals of}{274–275, 274t}
\idxentry{2}{定性资料分析软件}{qualitative data analysis software}{275–276}
\idxentry{2}{撰写与分享结果}{writing and sharing results}{277}
\idxentry{1}{数据类型与值域}{types, and domains}{343}
\idxentry{0}{数据与安全监察委员会（DSMB）}{Data and Safety Monitoring Board (DSMB)}{210}
\idxentry{0}{资料收集}{Data collection}{}
\idxentry{1}{为研究目的收集的既有数据／标本}{existing data/specimens for research purposes}{291}
\idxentry{1}{方案的小幅修订}{minor protocol revisions}{329–330}
\idxentry{1}{在线表单}{online forms}{345}
\idxentry{1}{平台}{platforms}{327}
\idxentry{1}{方案的实质修订}{substantive protocol revisions}{330}
\idxentry{0}{数据字典}{Data dictionary}{343, 404g}
\idxentry{0}{数据录入}{Data entry}{}
\idxentry{1}{数据管理软件}{data management software}{346–347, 347t}
\idxentry{1}{下拉列表}{drop-down list}{345, 346f}
\idxentry{1}{电子数据采集}{electronic data capture}{345, 346f}
\idxentry{1}{导入测量值}{importing measurements}{345–346}
\idxentry{1}{实验室结果}{laboratory results}{345–346}
\idxentry{1}{选项组}{option group}{345, 346f}
\idxentry{1}{回答}{responses}{}
\sourcepage{755}
\idxentry{2}{穷尽全部可能}{collectively exhaustive}{345}
\idxentry{2}{相互排斥}{mutually exclusive}{345}
\idxentry{1}{系统}{system}{340, 404g}
\idxentry{0}{数据管理}{Data management}{}
\idxentry{1}{数据分析}{analysis of data}{349}
\idxentry{1}{保密与安全}{confidentiality and security}{349}
\idxentry{1}{数据录入}{data entry}{}
\idxentry{2}{数据管理软件}{data management software}{346–347, 347t}
\idxentry{2}{电子数据采集}{electronic data capture}{345, 346f}
\idxentry{2}{导入测量与实验室结果}{importing measurements and laboratory results}{345–346}
\idxentry{1}{数据表}{data tables}{340–342, 341–343f}
\idxentry{2}{通用数据元素}{common data elements}{344–345}
\idxentry{2}{数据字典、类型与值域}{data dictionaries, data types, and domains}{343}
\idxentry{2}{变量名与编码约定}{variable names and coding conventions}{343–344, 344f}
\idxentry{1}{习题}{exercises for}{351–352}
\idxentry{1}{提取数据}{extracting data}{348f}
\idxentry{2}{识别并纠正数据错误}{identifying and correcting errors in data}{348–349}
\idxentry{1}{质量控制}{quality control for}{334–335, 335t}
\idxentry{1}{软件}{software}{346–347, 347t, 404g}
\idxentry{0}{数据挖掘}{Data-mining}{61}
\idxentry{0}{数据表}{Data table}{340–342, 341–343f, 404g}
\idxentry{1}{用于队列研究}{for cohort study}{341f}
\idxentry{1}{通用数据元素}{common data elements}{344–345}
\idxentry{1}{数据字典、类型与值域}{data dictionaries, data types, and domains}{343}
\idxentry{1}{简化形式}{simplified}{340, 341f}
\idxentry{1}{单个数据表}{single}{340}
\idxentry{1}{变量名与编码约定}{variable names and coding conventions}{343–344, 344f}
\idxentry{0}{数据使用协议}{Data use agreement}{299, 404g}
\idxentry{0}{数据库设计}{Database design}{328}
\idxentry{0}{数据库管理软件}{Database management software}{340}
\idxentry{0}{数据库查询}{Database query}{348, 348f, 405g}
\idxentry{0}{入组日期}{Date of entry}{145}
\idxentry{0}{检查对结局作用研究中的决策分析}{Decision analysis, studies of effect of testing on outcomes}{256}
\idxentry{0}{缺乏决策能力}{Decision-making capacity, lacking}{101}
\idxentry{0}{专用设备}{Dedicated equipment}{325–326}
\idxentry{0}{演绎}{Deduction}{266, 405g}
\idxentry{0}{去标识化数据}{Deidentified data}{298–299, 405g}
\idxentry{0}{因变量}{Dependent variable}{5。另见结局变量（Outcome variable）}
\idxentry{0}{描述性研究}{Descriptive study}{4, 53, 116, 268, 268t, 405g}
\idxentry{1}{指南}{guides}{273}
\idxentry{1}{样本量估计方法}{sample size techniques for}{75–76}
\sourcepage{756}
\idxentry{2}{连续变量}{continuous variables}{76, 91}
\idxentry{2}{二分类变量}{dichotomous variables}{76–77, 92}
\idxentry{1}{撰写与分享结果}{writing and sharing results}{277}
\idxentry{0}{研究设计。参见研究设计}{Design of study. See Study design}{}
\idxentry{0}{设计阶段}{Design phase}{}
\idxentry{1}{混杂因素的处理}{coping with confounders in}{164–165, 165t}
\idxentry{2}{匹配}{matching}{166–167}
\idxentry{2}{机会性研究}{opportunistic studies}{168–171}
\idxentry{2}{限定}{specification}{165–166}
\idxentry{1}{减少偏倚所致错误}{minimizing errors due to bias}{161–163}
\idxentry{0}{设备数据}{Device data}{299–300}
\idxentry{0}{诊断试验研究}{Diagnostic test study}{245, 405g}
\idxentry{1}{设计或分析的陷阱}{pitfalls in design or analysis of}{}
\idxentry{2}{删除临界或无法解释的结果}{borderline or uninterpretable results, dropping}{257}
\idxentry{2}{差异性验证偏倚}{differential verification bias}{258}
\idxentry{2}{不适当排除}{exclusion, inappropriate}{257}
\idxentry{2}{样本量不足}{sample size, inadequate}{257}
\idxentry{2}{采用单一金标准}{single gold standard, application of}{257–258}
\idxentry{2}{验证偏倚}{verification bias}{257–258}
\idxentry{0}{诊断收益研究}{Diagnostic yield studies}{253}
\idxentry{0}{作为自报测量的日记}{Diaries, as self-reported measures}{307, 313–314}
\idxentry{0}{二分类变量}{Dichotomous variable}{40, 76–77, 92, 405g}
\idxentry{1}{与连续变量的比较}{continuous variables vs.}{78–79}
\idxentry{1}{连续变量描述性研究的样本量}{sample size for descriptive study of}{92}
\idxentry{0}{膳食纤维摄入}{Dietary fiber consumption}{190}
\idxentry{0}{双重差分}{Difference-in-differences}{170–171, 232, 405g}
\idxentry{0}{差异性测量偏倚}{Differential measurement bias}{143–145}
\idxentry{0}{差异性错分偏倚}{Differential misclassification bias}{144, 405g}
\idxentry{0}{差异性验证偏倚}{Differential verification bias}{258, 405g}
\idxentry{0}{数字传感器}{Digital sensors}{111}
\idxentry{0}{有向无环图}{Directed acyclic graph (DAG)}{182, 406g}
\idxentry{1}{碰撞点}{collider}{175, 175f}
\idxentry{1}{连接变量}{connecting variables}{184–186}
\idxentry{2}{混杂}{confounding}{185, 185f}
\idxentry{2}{循环图}{cyclic graph}{185, 185f}
\idxentry{2}{中介作用}{mediation}{184, 184f}
\idxentry{2}{共同效应}{shared effects}{184–185, 185f}
\idxentry{1}{例子}{example}{182, 188–190}
\idxentry{1}{因果推断}{inferences about causality}{186–187, 186–187t}
\idxentry{1}{局限}{limitations}{190}
\idxentry{1}{未连接全部变量}{not connecting variables}{182–184, 183–184f}
\sourcepage{757}
\idxentry{1}{与抽样的关系}{sampling and}{187}
\idxentry{1}{共有原因}{sharing cause}{188}
\idxentry{0}{披露的信息}{Disclosed information}{}
\idxentry{1}{向参与者披露}{to participants}{98–99}
\idxentry{2}{对信息的理解}{understanding}{99}
\idxentry{0}{离散变量}{Discrete variable}{40, 406g}
\idxentry{0}{区分度}{Discrimination}{245, 248, 406g}
\idxentry{0}{疾病严重程度范围}{Disease severity, spectrum of}{241–242}
\idxentry{0}{分散式试验}{Disseminated trials}{209}
\idxentry{0}{回答分布}{Distribution of responses}{47}
\idxentry{0}{在多样人群中评估测量}{Diverse populations, evaluating measures for}{318}
\idxentry{0}{剂量反应}{Dose-response}{178, 406g}
\idxentry{0}{双重问题}{Double-barreled questions}{310}
\idxentry{0}{双队列研究}{Double-cohort study}{126f, 127–128, 406g}
\idxentry{0}{双金标准偏倚。参见差异性验证偏倚}{Double gold standard bias. See Differential verification bias}{}
\idxentry{0}{正式实施前的演练}{Dress rehearsals}{329}
\idxentry{0}{退出者}{Dropout}{72, 406g}
\idxentry{0}{DSMB。参见数据与安全监察委员会}{DSMB. See Data and Safety Monitoring Board (DSMB)}{}
\idxentry{0}{握力计}{Dynamometer}{50}
\idxletter{E}
\idxentry{0}{电子邮件调查}{E-mailed surveys}{307}
\idxentry{0}{职业早期研究者}{Early stage investigators}{353}
\idxentry{0}{生态学谬误}{Ecologic fallacy}{302, 406g}
\idxentry{0}{社会与经济不利处境}{Economic disadvantages, social and}{104–105}
\idxentry{0}{eCRFs。参见电子病例报告表}{eCRFs. See Electronic case report forms (eCRFs)}{}
\idxentry{0}{果因关系}{Effect-cause}{158, 159t, 163–164, 163t, 406–407g}
\idxentry{0}{效应修饰}{Effect modification}{165, 200, 221, 406g}
\idxentry{1}{与混杂的关系}{confounding and}{191–192}
\idxentry{1}{演示用假想例子}{hypothetical examples to demonstrate}{191–192}
\idxentry{1}{与交互作用的关系}{and interaction}{171, 176}
\idxentry{0}{效应量}{Effect size}{56, 64, 65, 406g}
\idxentry{0}{实际效果}{Effectiveness}{4, 407g}
\idxentry{0}{干预疗效}{Efficacy of intervention}{196}
\idxentry{0}{效率}{Efficiency}{47, 407g}
\idxentry{0}{嵌入电子病历的随机试验}{EHR-embedded randomized trials}{303}
\idxentry{0}{电子病例报告表（eCRFs）}{Electronic case report forms (eCRFs)}{345}
\idxentry{0}{电子数据采集}{Electronic data capture}{345, 346f}
\idxentry{0}{电子日志}{Electronic logs}{307, 313–314}
\idxentry{0}{流行病学家}{Epidemiologist}{59, 138, 188, 407g}
\idxentry{0}{流行病学}{Epidemiology}{120, 407g}
\idxentry{0}{临床均衡}{Equipoise}{108, 229, 407g}
\sourcepage{758}
\idxentry{0}{等效性}{Equivalence}{221–224, 223f}
\idxentry{0}{等效性研究}{Equivalence study}{74–75, 407g}
\idxentry{0}{研究中的错误}{Errors of research}{8–10, 9f, 10t}
\idxentry{0}{必需的专业知识}{Essential expertise}{20}
\idxentry{0}{伦理问题}{Ethical issues}{}
\idxentry{1}{伦理原则}{ethical principles}{95–96}
\idxentry{1}{习题}{exercises for}{114}
\idxentry{1}{人体研究联邦法规}{federal regulations for research on human subjects}{96–97}
\idxentry{2}{知情同意的例外}{exceptions to informed consent}{99–101}
\idxentry{2}{知情且自愿的同意}{informed and voluntary consent}{98–99}
\idxentry{2}{伦理审查委员会批准}{institutional review board approval}{97–98, 98t}
\idxentry{2}{尽量减少风险}{minimizing risks}{101–102}
\idxentry{2}{易受伤害研究参与者的保护}{protections for vulnerable research participants}{102–105}
\idxentry{1}{研究者的责任}{investigators, responsibilities of}{}
\idxentry{2}{署名}{authorship}{106–107}
\idxentry{2}{利益冲突}{conflicts of interest}{107–108}
\idxentry{2}{科研不端}{scientific misconduct}{105–106}
\idxentry{1}{其他问题}{other issues}{}
\idxentry{2}{研究的临床影响}{clinical impact of research}{111–112}
\idxentry{2}{向参与者付费}{payment to research participants}{111}
\idxentry{1}{临床研究法规的历史}{regulations on clinical research, history of}{95}
\idxentry{1}{特定研究类型的问题}{specific to certain types of research}{}
\idxentry{2}{使用人工智能和大数据的研究}{artificial intelligence and big data, research using}{110}
\idxentry{2}{既往收集标本和数据的研究}{previously collected specimens and data, research on}{109}
\idxentry{2}{随机临床试验}{randomized clinical trials}{108–109}
\idxentry{2}{使用传感器与移动健康设备的试验}{trials using sensors and mobile health devices}{110–111}
\idxentry{0}{伦理原则}{Ethical principles}{95–96}
\idxentry{0}{合乎伦理的研究}{Ethical research}{21–22}
\idxentry{0}{民族志}{Ethnography}{263, 270, 407g}
\idxentry{0}{Eureka研究平台}{Eureka Research Platform}{300}
\idxentry{0}{支持因果关系的依据}{Evidence favoring causality}{178–179}
\idxentry{0}{知情同意的例外}{Exceptions to informed consent}{99–101}
\idxentry{0}{排除标准}{Exclusion criteria}{5, 30, 201, 407g}
\idxentry{0}{免于伦理审查}{Exempt from IRB review}{97, 98t, 407g}
\idxentry{0}{既有数据／标本}{Existing data/specimens}{}
\idxentry{1}{创造性用途}{creative uses of}{}
\idxentry{2}{附属研究}{ancillary studies}{301–302}
\idxentry{2}{合并多来源数据}{combining data from multiple sources}{302}
\idxentry{2}{荟萃分析}{meta-analysis}{300–301}
\idxentry{2}{随机试验}{randomized trials}{302–303}
\idxentry{1}{习题}{exercises for}{306}
\sourcepage{759}
\idxentry{1}{来源}{sources}{}
\idxentry{2}{为研究目的收集资料}{data collection for research purposes}{291}
\idxentry{2}{医疗服务数据}{health care delivery data}{297–299}
\idxentry{2}{互联网与设备数据}{internet and device data}{299–300}
\idxentry{2}{全国调查}{national surveys}{291, 293–294}
\idxentry{2}{其他研究}{other research studies}{295–296}
\idxentry{2}{登记}{registries}{294–295}
\idxentry{2}{研究数据}{research data}{296}
\idxentry{1}{用途、优点与陷阱}{uses, advantages, and pitfalls of}{292–293t}
\idxentry{0}{伦理快速审查}{Expedited IRB review}{98, 98t, 407g}
\idxentry{0}{有经验的主要导师}{Experienced primary mentor}{23}
\idxentry{0}{实验}{Experiment}{53。另见临床试验（Clinical trial）}
\idxentry{1}{与分析性研究的关系}{analytic studies and}{65, 66t}
\idxentry{2}{卡方检验}{chi-squared test}{68–71, 87–89}
\idxentry{2}{相关系数}{correlation coefficient}{71–72, 90}
\idxentry{2}{t检验}{t test}{66–68, 86, 93}
\idxentry{0}{专业知识}{Expertise}{18}
\idxentry{0}{探索性研究}{Exploratory study}{267–268, 268t, 407g}
\idxentry{1}{指南}{guides}{273}
\idxentry{1}{撰写与分享结果}{writing and sharing results}{277}
\idxentry{0}{暴露}{Exposure}{5, 407–408g}
\idxentry{0}{外部对照}{External controls}{127–128}
\idxentry{0}{外部效度}{External validity}{6, 408g}
\idxletter{F}
\idxentry{0}{捏造}{Fabrication}{106, 408g}
\idxentry{0}{表面效度}{Face validity}{46, 317t, 408g}
\idxentry{0}{析因试验}{Factorial trial}{221, 222f, 408g}
\idxentry{0}{假阴性结果}{False-negative result}{246, 408g}
\idxentry{0}{假阳性结果}{False-positive result}{246, 408g}
\idxentry{0}{篡改}{Falsification}{106, 408g}
\idxentry{0}{证伪检验}{Falsification tests}{162, 231, 408g}
\idxentry{0}{可行性}{Feasibility}{33, 240}
\idxentry{1}{与研究问题的关系}{research question and}{19–21}
\idxentry{1}{可行性研究}{study of}{254}
\idxentry{0}{粪便隐血试验（FOBT）}{Fecal occult blood testing (FOBT)}{190}
\idxentry{0}{联邦健康隐私规则}{Federal Health Privacy Rule}{102}
\idxentry{0}{人体研究联邦法规}{Federal regulations for research on human subjects}{96–97}
\idxentry{1}{知情同意的例外}{exceptions to informed consent}{99–101}
\idxentry{1}{知情且自愿的同意}{informed and voluntary consent}{98–99}
\idxentry{1}{伦理审查委员会批准}{institutional review board approval}{97–98, 98t}
\idxentry{1}{尽量减少风险}{minimizing risks}{101–102}
\sourcepage{760}
\idxentry{1}{易受伤害研究参与者的保护}{protections for vulnerable research participants}{102–105}
\idxentry{0}{研究现场管理}{Field sites management}{272}
\idxentry{0}{字段}{Field}{340, 408g}
\idxentry{0}{经济利益冲突}{Financial conflicts of interest}{107, 408–409g}
\idxentry{0}{良好研究问题与计划的FINE标准}{“FINE” criteria for good research question and study plan}{19–22, 20t}
\idxentry{0}{撒网式探索}{Fishing expedition}{61}
\idxentry{0}{固定样本量}{Fixed sample size}{77–78}
\idxentry{0}{平面文件}{Flat file}{340}
\idxentry{0}{FOBT。参见粪便隐血试验}{FOBT. See Fecal occult blood testing (FOBT)}{}
\idxentry{0}{焦点小组}{Focus group}{263, 272, 314, 409g}
\idxentry{0}{随访}{Follow-up}{207–209, 207t}
\idxentry{1}{尽可能完成随访}{maximizing}{132–133, 133t}
\idxentry{0}{外键}{Foreign key}{341, 409g}
\idxentry{0}{表格设计}{Forms development}{327–328}
\idxentry{0}{基金会与专业学会}{Foundations and professional societies}{364–365, 365f}
\idxentry{0}{框架分析}{Framework analysis}{276, 409g}
\idxentry{0}{Framingham心脏研究}{Framingham Heart Study}{27, 296}
\idxentry{0}{伪造数据}{Fraudulent data}{334}
\idxentry{0}{频数匹配}{Frequency matching}{166}
\idxentry{0}{频率学派分析}{Frequentist analysis}{224, 235, 409g}
\idxentry{0}{完全序贯设计}{“Fully sequential” design}{237}
\idxentry{0}{获得资助的可能性}{Fundability}{20}
\idxentry{0}{基础社区参与}{Fundamental community engagement}{286}
\idxentry{0}{研究资助支持}{Funding support for research}{361–362}
\idxentry{1}{基金会与专业学会资助}{foundations and professional societies, grants from}{364–365, 365f}
\idxentry{1}{产业资助}{industry funding}{365–366}
\idxentry{1}{机构内部研究经费}{intramural research funds}{366}
\idxentry{1}{NIH资助与合同}{NIH grants and contracts}{362–364, 363t, 364f}
\idxentry{1}{产业研究支持}{research support from industry}{365–366}
\idxletter{G}
\idxentry{0}{GCP。参见药物临床试验质量管理规范}{GCP. See Good Clinical Practice (GCP)}{}
\idxentry{0}{可推广性}{Generalizability}{6, 283。另见外部效度（External validity）}
\idxentry{0}{用于产生假设}{Generating hypothesis, usefulness for}{140}
\idxentry{0}{基于基因组的预测模型}{Genome-based prediction model}{106}
\idxentry{0}{地域标准}{Geographic criteria}{28}
\idxentry{0}{幽灵作者}{Ghost authors}{107, 409g}
\idxentry{0}{金标准}{Gold standard}{44, 242, 409g}
\idxentry{1}{其中的不完善之处}{imperfections in}{242}
\idxentry{1}{单个数据表}{single}{257–258}
\idxentry{0}{药物临床试验质量管理规范（GCP）}{Good Clinical Practice (GCP)}{206, 331–332, 331t, 409g}
\idxentry{0}{良好申请书的特征}{Good proposal, characteristics of}{361}
\sourcepage{761}
\idxentry{0}{良好的研究假设}{Good research hypothesis}{}
\idxentry{1}{事先假设与事后假设}{advance vs. after-the-fact hypothesis}{54}
\idxentry{1}{简单与复杂假设}{simple vs. complex}{53}
\idxentry{1}{具体与含糊假设}{specific vs. vague}{54}
\idxentry{0}{良好的研究问题}{Good research question}{19–22}
\idxentry{0}{研究资助}{Grants}{}
\idxentry{1}{资助与合同}{and contracts}{362–364, 363t, 364f}
\idxentry{1}{基金会和专业学会资助}{from foundations and professional societies}{364–365, 365f}
\idxentry{1}{美国国立卫生研究院}{National Institutes of Health}{362–364, 363t, 364f}
\idxentry{0}{握力的操作定义}{Grip strength, operational definition of}{50}
\idxentry{0}{扎根理论}{Grounded theory}{267, 409g}
\idxentry{0}{挂名作者}{Guest author}{107, 409g}
\idxletter{H}
\idxentry{0}{风险率}{Hazard rate}{409g}
\idxentry{0}{风险率比}{Hazard ratio}{74, 132, 250, 409g}
\idxentry{0}{医疗服务中的临床试验}{Health care, clinical trials in}{206}
\idxentry{0}{医疗服务数据}{Health care delivery data}{297–299}
\idxentry{0}{Health eHeart研究}{Health eHeart Study}{295}
\idxentry{0}{健康公平}{Health equity}{22}
\idxentry{0}{健康保险流通与责任法案}{Health Insurance Portability and Accountability Act (HIPAA)}{102, 349, 410g}
\idxentry{1}{授权}{authorization}{298}
\idxentry{0}{健康素养}{Health literacy}{311, 410g}
\idxentry{0}{健康使用者效应}{Healthy user effect}{297, 410g}
\idxentry{0}{心脏与雌激素／孕激素替代研究（HERS）}{Heart and Estrogen/Progestin Replacement Study (HERS)}{296}
\idxentry{0}{HERS。参见心脏与雌激素／孕激素替代研究}{HERS. See Heart and Estrogen/Progestin Replacement Study (HERS)}{}
\idxentry{0}{异质性}{Heterogeneity}{300, 410g}
\idxentry{0}{问题中的隐含假设}{Hidden assumptions, questions and}{310}
\idxentry{0}{最高后验密度区间}{Highest posterior density interval}{236, 410g}
\idxentry{0}{HIPAA。参见健康保险流通与责任法案}{HIPAA. See Health Insurance Portability and Accountability Act (HIPAA)}{}
\idxentry{0}{同质性}{Homogeneity}{410g}
\idxentry{0}{医院来源对照}{Hospital-based controls}{142, 410g}
\idxentry{0}{人体研究参与者}{Human subjects}{}
\idxentry{1}{定义}{defined}{96}
\idxentry{1}{研究}{research}{21–22。另见人体研究联邦法规（Federal regulations for research on human subjects）}
\idxentry{1}{申请书中的部分}{section of proposal}{360–361}
\idxentry{0}{假设}{Hypothesis}{6, 410g}
\idxentry{1}{良好研究假设的特征}{characteristics of good research hypothesis}{}
\idxentry{2}{事先假设与事后假设}{advance vs. after-the-fact hypothesis}{54}
\idxentry{2}{简单与复杂假设}{simple vs. complex}{53}
\idxentry{2}{具体与含糊假设}{specific vs. vague}{54}
\sourcepage{762}
\idxentry{1}{假设的产生}{generation}{61}
\idxentry{1}{零假设与备择假设}{null and alternative hypothesis}{54–55}
\idxletter{I}
\idxentry{0}{不死时间偏倚}{Immortal time bias}{129, 410g}
\idxentry{0}{实施研究}{Implementing study}{8, 8f}
\idxentry{0}{研究问题的重要性}{Importance, research question}{21}
\idxentry{0}{不精密的数据}{Imprecise data}{334}
\idxentry{0}{不准确且不精密的数据}{Inaccurate and imprecise data}{334}
\idxentry{0}{累积发病率}{Incidence}{116, 117, 117t, 410g}
\idxentry{0}{发病密度病例对照研究}{Incidence-density case-control study}{138}
\idxentry{0}{发病密度巢式病例对照研究}{Incidence-density nested case-control study}{146–148, 147f}
\idxentry{1}{阿片过量死亡}{opiate overdose deaths}{147–148}
\idxentry{1}{优点}{strengths}{149}
\idxentry{1}{缺点}{weaknesses}{149}
\idxentry{0}{发病密度抽样}{Incidence-density sampling}{411g}
\idxentry{0}{发病密度}{Incidence rate}{131, 410g}
\idxentry{0}{纳入}{Inclusion}{5}
\idxentry{0}{纳入标准}{Inclusion criteria}{29, 200, 411g}
\idxentry{0}{不协调的问题}{Incongruous question}{311}
\idxentry{0}{纳入偏倚}{Incorporation bias}{242, 411g}
\idxentry{0}{独立}{Independent}{171, 411g}
\idxentry{0}{自变量}{Independent variable}{5。另见预测变量（Predictor variable）}
\idxentry{0}{待评价检验}{Index test}{242, 411g}
\idxentry{0}{归纳}{Induction}{266, 411g}
\idxentry{0}{产业资助}{Industry funding}{365–366}
\idxentry{0}{推断}{Inference}{5, 411g}
\idxentry{1}{关于因果关系的推断}{about causality}{186–187, 186–187t}
\idxentry{0}{知情同意}{Informed consent}{30, 95, 98–99, 411g}
\idxentry{1}{知情同意书}{consent forms}{99}
\idxentry{1}{向参与者披露信息}{disclosure of information to participants}{98–99}
\idxentry{1}{例外}{exceptions to}{99–101}
\idxentry{2}{缺乏决策能力}{decision-making capacity, lacking}{101}
\idxentry{2}{去标识化的剩余标本和数据}{leftover de-identified specimens and data}{100}
\idxentry{2}{反对意见}{objections to}{101}
\idxentry{2}{理由}{rationale for}{101}
\idxentry{1}{参与者对披露信息的理解}{participants’ understanding of disclosed information}{99}
\idxentry{1}{同意的自愿性质}{voluntary nature of consent}{99}
\idxentry{1}{豁免}{waiver of}{100, 100t}
\idxentry{0}{研究策略的创新部分}{Innovations section, of research strategy}{358}
\idxentry{0}{伦理审查委员会}{Institutional review board (IRB)}{95, 411g}
\idxentry{1}{批准}{approval}{97–98, 98t, 327}
\sourcepage{763}
\idxentry{2}{例外}{exceptions to}{97–98, 98t}
\idxentry{0}{仪器偏倚}{Instrument bias}{44, 411g}
\idxentry{0}{仪器变异}{Instrument variability}{42, 411g}
\idxentry{0}{工具变量}{Instrumental variable}{169, 412g}
\idxentry{0}{样本量估计中的信息不足}{Insufficient information, sample size estimation and}{82–83}
\idxentry{0}{预定样本}{Intended sample}{27, 412g}
\idxentry{0}{意向治疗分析}{Intention-to-treat analysis}{211–212, 213f, 412g}
\idxentry{0}{交互作用}{Interaction}{412g。另见效应修饰：与交互作用的关系（Effect modification, and interaction）}
\idxentry{0}{交互项}{Interaction term}{173, 412g}
\idxentry{0}{编码者间信度}{Intercoder reliability}{276, 412g}
\idxentry{0}{期中监测}{Interim monitoring}{210, 412g}
\idxentry{0}{中间标志物}{Intermediate marker}{198, 412g}
\idxentry{0}{内部一致性}{Internal consistency}{312, 412g}
\idxentry{0}{内部验证}{Internal validation}{253, 412–413g}
\idxentry{0}{内部效度}{Internal validity}{6, 413g}
\idxentry{0}{互联网和设备数据}{Internet and device data}{299–300}
\idxentry{0}{观察者间一致性}{Interobserver agreement}{261, 261t}
\idxentry{0}{观察者间可重复性}{Interobserver reproducibility}{243}
\idxentry{0}{解释}{Interpretation}{263, 413g}
\idxentry{0}{中断时间序列设计}{Interrupted time series design}{170, 232, 232f, 413g}
\idxentry{0}{干预}{Intervention}{3, 5, 21, 196, 413g}
\idxentry{1}{与对照条件的关系}{and control conditions}{197f}
\idxentry{2}{对照选择}{choice of control}{197–198}
\idxentry{2}{干预选择}{choice of intervention}{196–197}
\idxentry{2}{共同干预}{cointerventions}{198}
\idxentry{0}{访谈}{Interview}{272}
\idxentry{1}{测量实施}{administering measures}{320}
\idxentry{1}{指南}{guides}{273}
\idxentry{0}{访谈员施测的测量}{Interviewer-administered measures}{319}
\idxentry{0}{机构内部研究经费}{Intramural research funds}{366}
\idxentry{0}{观察者内可重复性}{Intraobserver reproducibility}{243}
\idxentry{0}{研究者的责任}{Investigators, responsibilities of}{}
\idxentry{1}{署名}{authorship}{106–107}
\idxentry{1}{利益冲突}{conflicts of interest}{107–108}
\idxentry{1}{科研不端}{scientific misconduct}{105–106}
\idxentry{0}{智商评分}{IQ score}{343}
\idxentry{0}{IRB。参见伦理审查委员会}{IRB. See Institutional review board (IRB)}{}
\idxentry{0}{迭代过程}{Iterative process}{22}
\idxletter{J}
\idxentry{0}{公正原则}{Justice}{96, 413g}
\idxletter{K}
\sourcepage{764}
\idxentry{0}{K类资助}{K awards}{362}
\idxentry{0}{K折交叉验证}{K-fold cross validation}{252, 413g}
\idxentry{0}{kappa}{Kappa}{42, 244, 413g}
\idxentry{1}{用于测量观察者间一致性的计算}{calculation to measure interobserver agreement}{261, 261t}
\idxletter{L}
\idxentry{0}{实验室程序的质量控制}{Laboratory procedures, quality control for}{333–334}
\idxentry{0}{从实验室到临床研究}{Laboratory to clinical research}{10–11}
\idxentry{0}{滞后变量}{Lagged variables}{146}
\idxentry{0}{领导力}{Leadership}{325}
\idxentry{0}{统计学显著性水平}{Level of statistical significance}{57, 64。另见\ensuremath{\alpha}水平（Alpha）}
\idxentry{0}{似然比（LR）}{Likelihood ratio (LR)}{76, 240, 249–250, 249t, 413g}
\idxentry{0}{李克特量表}{Likert scale}{308, 413g}
\idxentry{0}{有限数据集}{Limited data set}{299, 413–414g}
\idxentry{0}{地方能力}{Local capacity}{283–284}
\idxentry{0}{地方知识}{Local knowledge}{282–283, 283t}
\idxentry{0}{地方问题}{Local questions}{282–283, 283t}
\idxentry{0}{基于位置的数据}{Location-based data}{302}
\idxentry{0}{logistic回归模型}{Logistic regression model}{74, 132, 414g}
\idxentry{0}{纵向}{Longitudinal}{116, 414g}
\idxentry{0}{LRs。参见似然比}{LRs. See Likelihood ratios (LRs)}{}
\idxletter{M}
\idxentry{0}{机器学习算法}{Machine learning algorithm}{110}
\idxentry{0}{疟疾研究}{Malaria research}{111}
\idxentry{0}{边缘合计}{Marginals}{261, 414g}
\idxentry{0}{设盲。参见盲法}{Masking. See Blinding}{}
\idxentry{0}{匹配}{Matching}{74, 143, 165, 165t, 166, 178, 414g}
\idxentry{1}{优点}{advantages of}{166–167}
\idxentry{1}{病例对照研究中的匹配}{in case-control studies}{166}
\idxentry{1}{缺点}{disadvantages of}{167}
\idxentry{0}{成熟效应}{Maturation effects}{231, 414g}
\idxentry{0}{均值}{Mean}{40t, 42, 414g}
\idxentry{0}{使用多条目量表测量抽象变量}{Measure abstract variables, multi-item scales to}{311}
\idxentry{0}{测量}{Measurement}{39, 39f}
\idxentry{1}{准确度}{accuracy}{43–45, 44f, 44t}
\idxentry{2}{提高准确度的策略}{strategies for enhancing}{45–46, 45t}
\idxentry{1}{偏倚}{bias}{161–162}
\idxentry{1}{误差}{error}{9, 143, 414g}
\idxentry{1}{习题}{exercises for}{51}
\idxentry{1}{握力的操作定义}{grip strength, operational definition of}{50}
\idxentry{1}{测量方法的其他特征}{other features of approaches}{47}
\idxentry{1}{精密度}{precision}{41–42, 44t}
\sourcepage{765}
\idxentry{2}{评估}{assessing}{42}
\idxentry{2}{提高准确度的策略}{strategies for enhancing}{42–43, 43t}
\idxentry{1}{量表}{scales}{40t}
\idxentry{2}{分类变量}{categorical variables}{40}
\idxentry{2}{选择}{choosing}{41}
\idxentry{2}{数值变量}{numerical variables}{40–41}
\idxentry{1}{保存的材料}{stored materials}{47–48, 48t}
\idxentry{1}{效度}{validity}{46–47}
\idxentry{0}{中位数}{Median}{40t, 414g}
\idxentry{0}{中介作用}{Mediation}{184, 184f, 414g}
\idxentry{0}{中介变量}{Mediator}{164, 414g}
\idxentry{0}{医疗支出面板调查（MEPS）}{Medical Expenditure Panel Survey (MEPS)}{294}
\idxentry{0}{医学检查}{Medical tests}{}
\idxentry{1}{准确度}{accuracy of}{}
\idxentry{2}{分析}{analysis}{247–250, 247t, 248f, 249t}
\idxentry{2}{设计}{designs}{245–247}
\idxentry{1}{临床预测模型}{clinical prediction models}{}
\idxentry{2}{建立模型}{developing}{251, 252f}
\idxentry{2}{验证}{validation}{252–253}
\idxentry{1}{检查结果对临床决策的影响}{effect of results on clinical decisions}{}
\idxentry{2}{临床决策的前后研究}{before/after studies of clinical decision-making}{253–254}
\idxentry{2}{诊断收益研究}{diagnostic yield studies}{253}
\idxentry{1}{检查对结局的影响}{effect of testing on outcomes}{}
\idxentry{2}{分析}{analysis}{256}
\idxentry{2}{结局选择}{choice of outcome}{255}
\idxentry{2}{临床试验}{clinical trials}{256}
\idxentry{2}{决策分析}{decision analysis}{256}
\idxentry{2}{观察性研究}{observational studies}{255–256}
\idxentry{1}{可行性、成本与风险}{feasibility, costs, and risks of}{}
\idxentry{2}{分析}{analysis}{254}
\idxentry{2}{设计}{designs}{254}
\idxentry{1}{研究的一般问题}{general issues for studies of}{240, 241t}
\idxentry{2}{盲法的重要性}{blinding, importance of}{242}
\idxentry{2}{过拟合的危险}{dangers of overfitting}{243}
\idxentry{2}{金标准的不完善}{gold standard, imperfections in}{242}
\idxentry{2}{疾病严重程度和检查结果范围}{spectrum of disease severity and of test results}{241–242}
\idxentry{2}{变异来源}{variability, sources of}{242}
\idxentry{1}{用\ensuremath{\kappa}测量观察者间一致性}{kappa to measure interobserver agreement}{261, 261t}
\idxentry{2}{习题}{exercises for}{262}
\idxentry{1}{诊断检查研究设计或分析的陷阱}{pitfalls in design or analysis of diagnostic test studies}{}
\idxentry{2}{删除临界或无法解释的结果}{borderline or uninterpretable results, dropping}{257}
\sourcepage{766}
\idxentry{2}{差异性验证偏倚}{differential verification bias}{258}
\idxentry{2}{不适当排除}{exclusion, inappropriate}{257}
\idxentry{2}{样本量不足}{sample size, inadequate}{257}
\idxentry{2}{采用单一金标准}{single gold standard, application of}{257–258}
\idxentry{2}{验证偏倚}{verification bias}{257–258}
\idxentry{1}{可重复性}{reproducibility}{243}
\idxentry{2}{分析}{analysis}{244–245}
\idxentry{2}{设计}{designs}{243–244}
\idxentry{1}{用途}{usefulness of}{240–243}
\idxentry{0}{医学／公共卫生}{Medical/public health}{18}
\idxentry{0}{备忘录}{Memo}{2710, 414g}
\idxentry{0}{谅解备忘录}{Memorandum of Understanding (MOU)}{288, 415g}
\idxentry{0}{孟德尔随机化}{Mendelian randomization}{168–169, 415g}
\idxentry{0}{导师的选择与合作}{Mentor, choosing and working with}{22–23}
\idxentry{0}{MEPS。参见医疗支出面板调查}{MEPS. See Medical Expenditure Panel Survey (MEPS)}{}
\idxentry{0}{荟萃分析}{Meta-analysis}{31, 300–301, 415g}
\idxentry{0}{荟萃回归}{Meta-regression}{300, 415g}
\idxentry{0}{元数据}{Metadata}{343}
\idxentry{0}{微生物组的组成}{Microbiome, components of}{189}
\idxentry{0}{最小临床重要差异}{Minimal clinically important difference}{317t, 415g}
\idxentry{0}{最小风险的定义}{Minimal risk, defined}{98}
\idxentry{0}{最低必需数据}{Minimum necessary data}{298, 415g}
\idxentry{0}{错分}{Misclassification}{125, 415g}
\idxentry{0}{缺失数据}{Missing data}{334, 415g}
\idxentry{0}{混合方法研究}{Mixed methods research}{358}
\idxentry{0}{移动健康设备}{Mobile health devices}{110–111}
\idxentry{0}{建模的简化例子}{Modeling, simplified example of}{193–194}
\idxentry{0}{MORE试验。参见雷洛昔芬多结局评价试验}{MORE trial. See Multiple Outcomes of Raloxifene Evaluation (MORE) trial}{}
\idxentry{0}{MOU。参见谅解备忘录}{MOU. See Memorandum of Understanding (MOU)}{}
\idxentry{0}{多条目工具的格式设计}{Multi-item instruments, formatting of}{312–313, 315}
\idxentry{0}{多条目量表}{Multi-item scales}{311–312}
\idxentry{0}{多队列研究}{Multiple cohort study}{127–128, 166, 415g}
\idxentry{0}{多个假设}{Multiple hypothesis}{60–62, 199}
\idxentry{1}{检验}{testing}{54, 210, 415g}
\idxentry{1}{互不相关}{unrelated}{62}
\idxentry{0}{雷洛昔芬多结局评价（MORE）试验}{Multiple Outcomes of Raloxifene Evaluation (MORE) trial}{330}
\idxentry{0}{乘法与加法交互作用}{Multiplicative interactions, additive vs.}{177}
\idxentry{0}{多变量调整}{Multivariable adjustment}{74}
\idxentry{0}{多变量分析}{Multivariable analysis}{172t, 173, 251, 416g}
\idxentry{0}{问卷中相互排斥的回答}{Mutually exclusive responses in questionnaire}{345}
